%\documentclass[12pt,oneside]{amsart}
\documentclass{scrartcl}

%\documentclass{article}

\usepackage{amsthm,amsmath,amssymb}
\usepackage[numbers]{natbib}

%\usepackage[colorlinks]{hyperref}
\usepackage{hyperref}
\hypersetup{
    colorlinks,%
    citecolor=black,%
    filecolor=black,%
    linkcolor=black,%
    urlcolor=black
}
\usepackage{hypernat}
\usepackage[top=2.5cm, bottom=2.5cm, left=2.8cm,
right=2.8cm]{geometry}
\usepackage{float}

%\setlength{\parskip}{0pt}
%\setlength{\parsep}{0pt}
%\setlength{\headsep}{0pt}
%\setlength{\topskip}{0pt}
%\setlength{\topmargin}{0pt}
%\setlength{\topsep}{0pt}
%\setlength{\partopsep}{0pt}

\linespread{1}

%\usepackage{marvosym}

%\textwidth = 6.5 in \textheight = 9.2 in \oddsidemargin = 0.0 in
%\evensidemargin = 0.0 in \topmargin = -0.0 in \headheight = 0.0 in
%\headsep = 0.0in \parskip = 0.0in \parindent = 0.0in
%\def\baselinestretch{0.871}


\newtheorem{theorem}{Theorem}[section]
\newtheorem{proposition}[theorem]{Proposition}
\newtheorem{problem}[theorem]{Problem}
\newtheorem{fact}[theorem]{Fact}
\newtheorem{lemma}[theorem]{Lemma}
\newtheorem{defn}[theorem]{Definition}
\newtheorem{corollary}[theorem]{Corollary}
\newtheorem{remark}{Remark}[section]
\newtheorem{example}[theorem]{Example}

\newcommand{\E}{{\mathbb E}}
\newcommand{\se}{{\mathcal{E}}}
\newcommand{\R}{{\mathbb R}}
\renewcommand{\P}{{\mathbb P}}
\newcommand{\ac}{{\mathcal{A}}}
\newcommand{\A}{{\mathcal{A}}}
\newcommand{\B}{{\mathcal{B}}}
\newcommand{\C}{{\mathcal{C}}}
\newcommand{\M}{{\mathcal{M}}}
\newcommand{\Mf}{{\mathfrak{M}}}
\newcommand{\T}{{\mathcal{T}}}
\newcommand{\Z}{{\mathcal{Z}}}
\newcommand{\K}{{\mathcal{K}}}
\newcommand{\lc}{{\mathcal{L}}}
\newcommand{\Ps}{{\mathcal{P}}}
\newcommand{\G}{{\mathcal{G}}}
\newcommand{\pa}{{\mathring{p}}}
\newcommand{\F}{{\cal F}}
\newcommand{\samp}{{\mathcal{X}}}
\newcommand{\X}{{\mathcal{X}}}
\newcommand{\hi}{{\hat{\Phi}}}
\newcommand{\data}{{\text{data}}}

\newcounter{rcnt}[section]
\renewcommand{\thercnt}{(\roman{rcnt})}

\def\qt#1{\qquad\text{#1}}

\def\argmin{\mathop{\rm argmin}}
\def\argmax{\mathop{\rm argmax}}


\setlength{\parskip}{1.4 \medskipamount}
%\sloppy
%\linespread{1.3}

\begin{document}

\title{STAT 238 - Bayesian Statistics \\
 Homework Two} 
\subtitle{Spring 2026, UC Berkeley}
\author{Due by 11:59 pm on 27 February 2026 \\ Total Points = 80} 

\date{}
\maketitle



% \tableofcontents

\begin{enumerate}

\item Suppose we are interested in the survival rate for a high-risk
  operation in a new hospital. Experience in other hospitals indicates
  that the probability of survival for a patient is expected to be
  around 0.9 and it would be fairly surprising if it were less than
  0.8 and more than 0.97. 
  \begin{enumerate}
  \item If you were to use the $\text{Beta}$ distribution to represent the
    opinion above, what values of $\alpha$ and $\beta$ would you use,
    and why? (\textbf{3 points}). 
  \end{enumerate}
  Suppose now that $n = 10$ patients are operated on in this new
  hospital and all of them survived.
  \begin{enumerate}
  \item What is the current posterior distribution of the survival
    rate? (\textbf{2 points})
  \item  What is the probability that the next patient will survive
    the operation? (\textbf{3 points})
  \item What is the probability that there will be 2 or more deaths in
    the next 20 operations? (\textbf{4 points}) 
  \end{enumerate}

\item Suppose we have a coin that we know is strongly biased either
  usually to come up heads or usually to come up tails but we don't
  know which.
  \begin{enumerate}
  \item Express this information in the form of a $\text{Beta}(\alpha,
    \beta)$ prior. What values of $\alpha$ and $\beta$ would you use
    and why? (\textbf{2 points})
   \item We now flip the coin 5 times and it comes up heads in 4 of
     the 5 flips. Based on your  prior from the previous part, what is
     the posterior distribution? (\textbf{2  points})
   \end{enumerate}

 \item Consider the Bayesian model:
   \begin{align*}
     &  \text{Prior}: \theta \sim \text{uniform}(0, 1) \\
     & \text{Likelihood}: X_i \overset{\text{i.i.d}}{\sim}
       \text{Bernoulli}(\theta). 
   \end{align*}
   Prove that
   \begin{align*}
     \P \left\{X_{n+1} = 1 \mid X_1= x_1, \dots, X_n = x_n \right\} =
     \frac{1 + \sum_{i=1}^n x_i}{n+2}. 
   \end{align*}
   This is called the Laplace rule of succession (see
   \url{https://en.wikipedia.org/wiki/Rule_of_succession}). Also
   compare with Problem 2(e) from Homework One. (\textbf{3 points}). 

 \item Consider the following Bayesian model:
   \begin{align*}
&   \text{Prior}: \log \frac{\theta}{1 - \theta} \sim N(\mu,
  \sigma^2) \\
& \text{Likelihood}: X \mid \theta \sim \text{Binomial}(n, \theta). 
   \end{align*}
   Assume that $n, \mu$ and $\sigma$ are fixed.
   \begin{enumerate}
   \item Show that the posterior of $\theta$ is given by (\textbf{2 points}):
\begin{equation*}
  f_{\theta | X = x}(\theta) = \frac{ \theta^x (1 - \theta)^{n-x} \frac{1}{\theta(1 - \theta)} \exp \left(-\frac{1}{2 \sigma^2}
    \left(\log \frac{\theta}{1 - \theta} - \mu \right)^2 \right) I\{0
  < \theta < 1\}}{\int_0^1 \theta^x (1 - \theta)^{n-x} \frac{1}{\theta(1 - \theta)} \exp \left(-\frac{1}{2 \sigma^2}
    \left(\log \frac{\theta}{1 - \theta} - \mu \right)^2 \right)
  d\theta }. 
\end{equation*}     
\item Suppose $x = 0$, $n = 1295$, $\mu = -9.27308950$ and $\sigma^{2} =
0.09259668$.  Calculate the posterior mean, posterior variance and a 95\%
  credible interval (an interval for which the posterior probability
  equals .95) by numerical discretization. (\textbf{6 points}).
 \item Solve part (b) using PyMC. Compare your answer in (b) with the
   answer obtained by PyMC. (\textbf{3 points}).  
\end{enumerate}

\item Consider the following Bayesian model:
  \begin{align*}
    & \text{Prior}: \theta \sim \text{Gamma}(\alpha, \beta) \text{ with
      density } \propto \theta^{\alpha - 1} \exp(-\beta \theta)
      I\{\theta > 0\} \\
    & \text{Likelihood}: Y \mid \theta \sim \text{Poisson}(k \theta)
  \end{align*}
  Treat $\alpha, \beta, k$ as fixed positive constants here.
  \begin{enumerate}
  \item Prove that the posterior of $\theta$ is given by: (\textbf{2 points})
    \begin{align*}
      \theta \mid Y = y \sim \text{Gamma}(\alpha + y, \beta + k). 
    \end{align*}
  \item Prove that the marginal distribution of $Y$ is given by:
    (\textbf{3 points})
    \begin{align*}
      \P\{Y = y\} = \frac{\Gamma(\alpha + y)}{\Gamma(\alpha) y!}
      \left(\frac{k}{\beta + k} \right)^y \left(\frac{\beta}{\beta +
      k} \right)^{\alpha} \qt{for $y = 0, 1, 2, \dots$}. 
    \end{align*}  
  \end{enumerate}
\item Suppose that the number of asthma deaths in a certain city in
   the United States with a population of 200000 in a single year can
   be modeled by the Poisson distribution with mean $2 \theta$ where
   $\theta$ represents the true underlying long-term ashtma mortality
   rate  in the city (measured in cases per 100000 persons per
   year). Suppose that in a particular year $3$ asthma deaths were
   observed in the city.
   \begin{enumerate}
   \item \textbf{Prior distribution for $\theta$}: Reviews of asthma
     mortality rates around the world suggest that mortality rates
     above 1.5 per 100000 people are rare in Western countries, with
     typical asthma mortality rates around 0.6 per 100000. Suggest a
     reasonable $\text{Gamma}(\alpha, \beta)$ prior for $\theta$ based
     on this information. Provide 
     reasoning for your answer. (\textbf{3 points})
   \item Compute the posterior distribution of $\theta$ given the data
     of observing $3$ deaths. (\textbf{3 points})
   \item Suppose now that ten years of data are obtained for this city
     (instead of just one) and 30 deaths are observed in the 10
     years. Calculate the posterior distribution of $\theta$ given
     this data (you can make assumptions such as the population
     remaining constant at 200000 over the 10 years etc.). (\textbf{3
       points})
   \item Comment on the similarities and differences between the two
     posteriors calculated above. (\textbf{2 points}). 
   \end{enumerate}

\item The kidney cancer dataset is analyzed in Section 2.7 of
  \citet{Gelman2013}. In that section, the authors use Poisson
  likelihoods with Gamma priors, whereas in class we worked with
  Binomial likelihoods and Beta priors. Here, we revisit the in-class
  analysis, replacing the Binomial–Beta model with a Poisson–Gamma
  model.

  As in class, let $X_i$ denote the total death count due to kidney
  cancer for county $i$ in the period 1980-89. Also let $n_i$ to be
  the average population for county $i$ in the same time period. Let
  $\theta_i$ be an unknown parameter representing the death rate for
  county $i$ in units of deaths per person per decade. 

  \begin{enumerate}
  \item Argue that $X_i \sim \text{Poisson}(n_i \theta_i)$. You can
    use the fact that $Bin(n, p)$ is well-approximated by
    $\text{Poisson}(np)$ provided $np^2$ is small (see e.g.,
    \url{https://en.wikipedia.org/wiki/Le_Cam\%27s_theorem}). (\textbf{2
      points})
  \item Consider the model:
    \begin{align*}
    &  \theta_i \overset{\text{i.i.d}}{\sim} \text{Gamma}(\alpha,
      \beta) \\
& X_i \mid \theta_i \sim \text{Poisson}(n_i \theta_i). 
    \end{align*}
    Treat $\alpha$ and $\beta$ as unknown parameters, and estimate
    them by maximizing the marginal likelihood. Use a numerical
    grid-based discretization approach to compute the maximum marginal
    likelihood estimates $\hat{\alpha}$ and $\hat{\beta}$ of $\alpha$
    and $\beta$. (\textbf{4 points}).
  \item Use $\hat{\alpha}, \hat{\beta}$ derived in (b) above to
    compute the posterior mean estimates of $\theta_i, i = 1, \dots,
    N$. Compare these estimates with those derived in class using the
    Beta-Binomial model. Identify the top 100 counties in terms of
    $\hat{\theta}_i$, and compare this list with that derived in class using the
    Beta-Binomial model. (\textbf{3 points})
  \item Consider the full Bayes model:
    \begin{equation}\label{fullBayes}
      \begin{split}
  & \alpha, \beta \sim \text{uninformative priors} \\
    &  \theta_i \overset{\text{i.i.d}}{\sim} \text{Gamma}(\alpha,
      \beta) \\
  & X_i \mid \theta_i \sim \text{Poisson}(n_i \theta_i).
    \end{split}
    \end{equation}
    for some suitable uninformative priors for $\alpha$ and
    $\beta$. What uninformative priors would you use and why?
  Without using any probabilistic programming library, implement
  Bayesian inference in this model, and compute posterior mean
  estimates of $\theta_1, \dots, \theta_N$. You can use the
  discretization based approach discussed in class for the
  Beta-Binomial model. (\textbf{5 points})
\item Use PyMC (or some other library) to obtain posterior samples for
  the model \eqref{fullBayes}, and use these samples to compute the
  posterior mean esimates of $\theta_1, \dots, \theta_N$. Compare the
  posterior means with those obtained in part (d) above. (\textbf{3
    points}) 
    \end{enumerate}
  \item Consider the following baseball dataset (from \citet[Chapter
    1]{Efron}) used in Lectures 10 and 11. The goal is to use the
    players' early season performance (as indicated by the second
    column) to predict their end of season performance (as indicated
    by the third column).

    Use the following model for this dataset:
\begin{align*}
 & a, b \overset{\text{i.i.d}}{\sim} \text{uniform}(-\infty, \infty) \\
 & \theta_i \mid a, b \sim \text{Beta}(a, b) \\
 & H_i \overset{\text{ind}}{\sim} \text{Bin}(45, \theta_i)
\end{align*}
where $H_i$ is the number of hits in the first $45$ at-bats for player
$i$. 

\begin{enumerate}
\item Using a numerical discretization approach, implement Bayesian
  inference in this model, and obtain posterior samples of $\theta_1,
  \dots, \theta_N$. Then compute the Bayesian estimates (posterior
  means) for each $\theta_i$. Compute the accuracy of these Bayesian
  estimates in the two metrics: sum of squares, and sum of absolute
  values. Compare the accuracies with the naive estimates $H_i/45$, as
  well as the accuracy of the James-Stein based estimates computed in
  Lecture 11. (\textbf{5 points}).
\item Using the posterior samples, compute 95\% credible intervals for
  each $\theta_i$. Compare these intervals with the actual values of
  $\theta_i$. How many of the actual $\theta_i$s lie in the intervals?
  (\textbf{3 points}). 
\end{enumerate}

\begin{table}[H]
\centering
\begin{tabular}{|l|p{0.25\linewidth}|c|}
\hline
\textbf{Player Name} & \textbf{Number of Hits in the first 45 At-Bats} & \textbf{EoSAverage} \\
\hline
Clemente     & 18 & .346 \\
F Robinson   & 17 & .298 \\
F Howard     & 16 & .276 \\
Johnstone    & 15 & .222 \\
Berry        & 14 & .273 \\
Spencer      & 14 & .270 \\
Kessinger    & 13 & .263 \\
L Alvarado   & 12 & .210 \\
Santo        & 11 & .269 \\
Swoboda      & 11 & .230 \\
Unser        & 10 & .264 \\
Williams     & 10 & .256 \\
Scott        & 10 & .303 \\
Petrocelli   & 10 & .264 \\
E Rodriguez  & 10 & .226 \\
Campaneris   & 9  & .286 \\
Munson       & 8  & .316 \\
Alvis        & 7  & .200 \\
\hline
\end{tabular}
\caption{Some Statistics of 18 Baseball Players from the 1970 Season}
\label{tab:player_stats}
\end{table}


\item In Lecture 11, we used the following fact:
  \begin{equation}\label{norm_lik}
    \begin{split}
  \theta \sim N(\mu, \tau^2) ~ \text{ and } ~ y \mid \theta \sim
  N(\theta, \sigma^2) &\implies \theta \mid y \sim N
  \left(\frac{y/\sigma^2 + \mu/\tau^2}{1/\sigma^2 + 1/\tau^2},
  \frac{1}{1/\sigma^2 + 1/\tau^2} \right) \\ &\text{ and } y \sim N(\mu,
                                               \sigma^2 + \tau^2).
  \end{split}                                               
\end{equation}
The goal of this problem is to prove a much more general form of this
result (this general form will be useful later on in the course).

Suppose $X$ and $Y$ are random
  vectors with $X \sim N(m_0, Q_0)$ and $Y | X = x \sim N(Bx, R)$. Let
  $X | Y = y \sim N(m_1, Q_1)$. 
  \begin{enumerate}
  \item Show that (\textbf{4 points})
    \begin{align*}
&      m_1 = (I + Q_0 B' R^{-1} B)^{-1} m_0 + Q_1 B' R^{-1} y \\ & Q_1
                                                                   =
                                                                   (I
                                                                   +
                                                                   Q_0
                                                                   B'
                                                                   R^{-1}
                                                                   B)^{-1}
                                                                   Q_0.  
    \end{align*}
  \item Show that (\textbf{3 points})
    \begin{align*}
      &m_1 = m_0 + Q_0 B' \left(B Q_0 B' + R \right)^{-1} (y - B m_0) \\
      &Q_1 = Q_0 - Q_0 B' \left(B Q_0 B' + R
      \right)^{-1} B Q_0. 
    \end{align*}
  \item Show that (\textbf{2 points})
    \begin{align*}
      Y \sim N(B m_0, BQ_0B^T + R). 
    \end{align*}
  \item Deduce \eqref{norm_lik} as a consequence of the results in
    parts (a)--(c). (\textbf{2 points})
  \end{enumerate}

\end{enumerate}


\begin{thebibliography}{99}

  
%\bibitem[Fisher(1934)]{Fisher1934}
%R.~A. Fisher,
%\newblock ``Probability, likelihood and quantity of information in the logic of
%uncertain inference,''
%\newblock \emph{Proceedings of the Royal Society of London. Series~A,
%Containing Papers of a Mathematical and Physical Character}, vol.~146,
%no.~856, pp.~1--8, 1934.

\bibitem[Efron(2010)]{Efron}
Efron, B. (2010)
\textit{Large-scale inference: empirical Bayes methods for estimation,
testing, and prediction}. Cambridge University Press, 2010.

\bibitem[Gelman et~al.(2013)]{Gelman2013}
Gelman, A., Carlin, J. B., Stern, H. S., Dunson, D. B., Vehtari, A., and Rubin, D. B. (2013)
\textit{Bayesian data analysis}. Third edition. Chapman \& Hall/CRC, 2013.

%     \bibitem[Jaynes(2003)]{Jaynes} Jaynes, E. T, (2003)
%  \textit{Probability theory: the logic of science}. Cambridge
%  University Press, 2003.

%  \bibitem[Sinai(1992)]{Sinai} Sinai, Y. G, (1992)
%  \textit{Probability theory: an introductory course}. Springer
%  Verlag, 1992.  

%\bibitem[{deGroot(1986)}]{DeGrootBlackwell}DeGroot, M. H. (1986)
%       A conversation with David Blackwell.
%\newblock \emph{Statistical Science}, \textbf{1}, 40--53.

%         \bibitem[Mosteller(1987)]{Mosteller} Mosteller, F, (1987)
%  \textit{Fifty challenging problems in probability with
%    solutions}. Courier Corporation, 1987.

%    \bibitem[MacKay(2003)]{MacKay} MacKay, D. J. C., (2003)
%  \textit{Information theory, inference, and learning
%  algorithms}. Cambridge University Press, 2003.

%\bibitem[{Lindley(2013)}]{Lindley}Lindley, D. V. (2013)
%   \textit{Understanding Uncertainty}. John Wiley and sons, 2013.


\end{thebibliography}


\end{document}

%%% Local Variables:
%%% mode: latex
%%% TeX-master: t
%%% End:
